% =====================================================================
%  Phase C --- Deterministic navigation
%  SCOPE: the free-space decomposition, the three-leg route, the brake.
% =====================================================================
\section{Deterministic Navigation}
\label{sec:phaseC_nav}

\subsection{Autonomous, but not exploratory}

This section states plainly what \phaseC{} navigation is, because the
distinction is easy to lose and a reader may reasonably assume that
removing SLAM removed autonomy.

The robot in \phaseC{} is fully autonomous in the control sense: it
receives a list of station labels and nothing else, and it decides its
own route, drives itself in closed loop, brakes on its own sensor
evidence, and corrects its own final position against what it sees.
There is no teleoperation and no trajectory replay.

What it does \emph{not} do is discover. It never builds a map, never
decides where to go next based on what it found, and never searches for
a path. It cannot, because the question does not arise: the map is a
catalogue that predates the robot (Sec.~\ref{sec:phaseC_catalogue}), and
the order names the stations explicitly. The autonomy that remains is
\emph{deterministic} rather than \emph{exploratory}, and the difference
is worth naming precisely rather than describing as ``less''.

\subsection{The free space, decomposed once}

The greenhouse free space is not a general planning problem. It is four
corridors joined at both ends, and \filename{agri/aisles.py} derives
that decomposition from the wall and gutter coordinates rather than
declaring it:

\begin{table}[!t]
\caption{The free bands, derived from the catalogue. Each is an interval
of the robot centre line in $y$; every station lies in exactly one.}
\label{tab:bands}
\centering\footnotesize
\begin{tabular}{@{}clc@{}}
\toprule
band & $y$ interval (m) & stations \\
\midrule
0 & $[-2.45,\; -1.40]$ & \station{P1,*R} \hfill 8 \\
1 & $[-1.00,\; -0.20]$ & \station{P1,*L}, \station{P2,*R} \hfill 16 \\
2 & $[+0.20,\; +1.00]$ & \station{P2,*L}, \station{P3,*R} \hfill 16 \\
3 & $[+1.40,\; +2.45]$ & \station{P3,*L} \hfill 8 \\
\bottomrule
\end{tabular}
\end{table}

At each end, past the \SI{8.0}{m} gutters, lies a \emph{headland} that
joins all four bands. A route is therefore at most three legs --- out to
a headland, across to the target band, in to the station --- and
choosing between the two headlands is a comparison of two numbers, not a
search:

\begin{lstlisting}[style=python]
def route(sx, sy, tx, ty):
    here, there = band(sy), band(ty)
    if here is not None and here == there:
        return [(tx, ty)]                 # one leg
    end = min(HEADLANDS,
              key=lambda e: abs(sx - e) + abs(e - tx))
    return [(end, sy), (end, ty), (tx, ty)]
\end{lstlisting}

\paragraph{Why no planner} A$^\star$ over an occupancy grid would spend
a millisecond rediscovering Table~\ref{tab:bands}, and would then be
free to discover something \emph{else} on a day when the grid was subtly
wrong. That is not a hypothetical: it is what \phaseA{} did in a
\SI{0.16}{m} clearance (Sec.~\ref{sec:phaseA_planning}). A closed-form
route cannot produce a novel answer, which in a corridor with
\SI{0.16}{m} of margin is a feature.

The cost of that choice is that the geometry must be right, so it is
verified exhaustively rather than argued: the suite sweeps
\textbf{all 2\,256 station-to-station routes}, plus every route out of
the dock, and asserts the robot's rectangle never comes within
\SI{0.08}{m} of a gutter or a wall on any leg of any of them. That runs
in about a second and needs nothing installed.

\subsection{The asymmetry that a symmetric model hides}

The two headlands are \textbf{not} at symmetric coordinates:
$x_{\text{west}} = \measured{-4.08}$~m and
$x_{\text{east}} = \measured{+4.87}$~m. The robot is not symmetric --- a
\SI{0.50}{m} sensor boom points $+x$ and \SI{0.29}{m} of chassis trails
behind --- so at the west end the sensor point leads and must stop short
of the gutters, while at the east end the sensor point goes deep and the
\emph{tail} must clear them. The footprint ends up centred in the
\SI{0.95}{m} gap either way, with \SI{0.08}{m} of margin at each end.

An earlier version used one number for both. Going east, that left the
chassis lying alongside the gutters while the robot strafed across the
greenhouse --- straight through them. \textbf{Nothing would have
reported it.} This robot carries no collision geometry, so driving
through a gutter produces no contact, no warning and no complaint, only
a recording that cannot be shown to anyone. It was caught by the
clearance sweep above, which is the entire reason that function exists.

\subsection{The brake, and what it must ignore}

The route is geometric; the lidar is not. During every leg the robot
tests each return against the corridor it is about to sweep. Two design
choices carry the section.

\paragraph{A corridor, not a cone} A $\pm35\degree$ cone was the first
version, and it stopped the robot dead in every aisle: at \SI{0.35}{m}
of lateral clearance, a gutter wall the robot is driving \emph{parallel}
to sits well inside a $35\degree$ cone, so the brake fired on the very
thing the route was designed to pass. A corridor asks the only question
that matters --- is this in the way? --- and a wall you are driving
alongside is not in the way.

The corridor's edges are the extreme projections of the robot's own four
corners onto the axis across the direction of travel. For a symmetric
robot a half-width would do; with a \SI{0.50}{m} boom in front and
\SI{0.29}{m} of chassis behind, a \emph{strafing} robot sweeps
\SI{0.79}{m} and not centred on the base.

\paragraph{The building is not an obstacle} The brake exists for things
that should not be there. The walls and the gutters should be there,
they are in the catalogue, and the sweep above has already proved every
route clears them. Braking for them is not caution; it is refusing to
move --- and it produced two failures that looked entirely different:

\begin{itemize}[leftmargin=*,itemsep=1pt,topsep=2pt]
  \item \emph{Every leg to a headland stopped.} The robot must back its
        tail to within \SI{0.08}{m} of the end wall to fit the
        \SI{0.95}{m} gap; the brake wanted \SI{0.35}{m} of clear floor
        and refused \SI{0.27}{m} short. The mission logged ``something
        is in the aisle''. Something was: the building.
  \item \emph{A leg that shifts sideways while advancing} ---
        \SI{0.10}{m} across while travelling \SI{0.66}{m} along, which
        is exactly what moving between the two stations of an inner
        aisle looks like --- tilts the swept corridor by $8.6\degree$
        and swings it onto the gutter the robot is driving \emph{beside}.
        Same wall, different geometry, same wrong answer.
\end{itemize}

Both are resolved by projecting each return into world coordinates and
discarding it when it falls within \SI{0.12}{m} of catalogued structure.
A return \emph{inside} the robot's own outline is discarded too: the
lidar sits within the chassis and sees the camera pedestal and the mast
at \SI{0.12}{m} and \SI{0.22}{m} ahead, which a raw range test brakes
for permanently, from the first metre.

The suite pins both behaviours from opposite sides --- a gutter being
driven past does not stop the robot, an object genuinely ahead does, and
no station stands on a point the brake would ignore --- so neither fix
can be undone by loosening the other.
